\section*{الفصل \ref{ch:b2:eruptive-styles} --- الثورانات البركانية}

\begin{solution}{exo:b2:eruptive-styles:1}
(أ) $\qty{1.3}{km^3} = \qty{1.3e9}{m^3}$، بين $10^9$ و$10^{10}$: VEI~5.
(ب) $\qty{2.9e8}{m^3}$: VEI~4. (ج) $\qty{2e6}{m^3}$، بين $10^6$ و$10^7$:
VEI~2. (د) يبدأ VEI~6 عند \qty{e10}{m^3}، فيجب أن تكون \omterm{def:b2:eruptive-styles:tephra}{التيفرا} الكلية
\qty{10}{km^3} على الأقل؛ و«أكثر من \qty{5}{km^3}» حدّ أدنى، لا تناقض.
\end{solution}

\begin{solution}{exo:b2:eruptive-styles:2}
(أ) رماد (أدقّ من \qty{2}{mm}). (ب) لابيلي (\qtyrange{2}{64}{mm}).
(ج) كتلة: أكبر من \qty{64}{mm}، زاويّة، صلبة حين قُذفت. (د)
\omterm{def:b2:eruptive-styles:tephra}{قنبلة بركانية}، قنبلة مغزلية: أكبر من \qty{64}{mm} وتشكّلت وهي منصهرة
في الطيران.
\end{solution}

\begin{solution}{exo:b2:eruptive-styles:3}
(أ) $0.304\times(10^6)^{0.241} = 0.304\times27.9 = \qty{8.5}{km}$؛
$0.304\times(10^8)^{0.241} = 0.304\times84.7 = \qty{25.8}{km}$.
(ب) $Q = (35/0.304)^{1/0.241} = 115.1^{4.149} = \qty{3.6e8}{kg/s}$، و$I =
8.55 + 3 = 11.6$. (ج) $Q \propto H^{4.15}$، فالارتفاع الأعلى من الحقيقي بنسبة \qty{20}{\percent}
يبالغ في تقدير $Q$ بعامل $1.2^{4.15} = 2.1$.
\end{solution}

\begin{solution}{exo:b2:eruptive-styles:4}
$V_{\mathrm{DRE}} = 0.29\times1000/2500 = \qty{0.116}{km^3}$؛ $m =
\num{0.29e9}\times1000 = \qty{2.9e11}{kg}$؛ $M = \log_{10}(\num{2.9e11}) - 7 = 11.46
- 7 = 4.46$.
\end{solution}

\begin{solution}{exo:b2:eruptive-styles:5}
$k = \ln(40/10)/(15 - 5) = 1.386/10 = \qty{0.139}{km^{-1}}$؛ $T_0 = 40\,e^{0.139\times5}
= 40\times2.00 = \qty{80}{cm}$؛ $\ln2/k = \qty{5.0}{km}$. ومع $T_0 =
\qty{8.0e-4}{km}$: $V = 2\times\num{8.0e-4}/0.139^2 = \qty{0.083}{km^3} =
\qty{8.3e7}{m^3}$، أي VEI~3.
\end{solution}

\begin{solution}{exo:b2:eruptive-styles:6}
(أ) $\int_0^x T_0e^{-ky}2y\,dy = 2T_0\big[-\tfrac{y}{k}e^{-ky} - \tfrac{1}{k^2}e^{-ky}\big]_0^x
= \tfrac{2T_0}{k^2}[1 - (1 + kx)e^{-kx}]$. (ب) بتكرار $u = \ln(2(1 + u))$ انطلاقًا من
$u = 1$: $1.386$، $1.564$، $1.635$، $1.662$، $1.672$، \dots، فيتقارب نحو
$u = 1.678$، التي يكون عندها $(1 + u)e^{-u} = 2.678\times0.1867 = 0.500$. (ج) $x_{1/2} =
1.678/0.139 = \qty{12.1}{km}$، داخل أبعد خط تساوي سمك ($\sqrt A =
\qty{15}{km}$)؛ ومع ذلك فالحجم الواقع خارج ذلك الخط هو $(1 + 2.09)e^{-2.09} =
\qty{38}{\percent}$ من الكلي. فنحو خُمسي الراسب يقعان حيث
لم يُقس شيء: والحجم يستند إلى الاستكمال.
\end{solution}

\begin{solution}{exo:b2:eruptive-styles:7}
(أ) $p_s = (5.0/0.444)^2 = \qty{127}{MPa}$، على عمق $127/0.025 = \qty{5.1}{km}$.
(ب) $x = (5.0 - 0.444\sqrt{10})/100 = (5.0 - 1.404)/100 = 0.0360$؛
$V_g = 0.0360\times8.314\times1123/(0.018015\times10^7) = \qty{1.86e-3}{m^3}$؛
$V_m = 0.964/2300 = \qty{4.19e-4}{m^3}$؛ $\phi = 1.86/(1.86 + 0.419) = 0.82$.
(ج) المصهور البازلتي أقل لزوجة من الريوليتي بآلاف إلى مليارات
المرات (\cref{ch:b2:magma-emplacement}): فتصعد فقاعاته،
وتندمج وتفلت، ويغادر الغاز عبر رغوة نفوذة أو عند
المنفذ، فتنزع \omterm{def:b2:magma-emplacement:magma}{الصهارة} غازها في نظام مفتوح بدلًا من أن تحتفظ به؛
وتفتّتها، حين يحدث، يقتصر على نافورة قرب المنفذ. والصعود الأبطأ
يترك وقتًا أطول لذلك.
\end{solution}

\begin{solution}{exo:b2:eruptive-styles:8}
(أ) $u_z(0) = 0.75\times10^7/(\pi\times(5000)^2) = \qty{0.095}{m} = \qty{95}{mm}$.
(ب) $z_0 = 3.8/0.766 = \qty{5.0}{km}$. (ج) $(1 + 1)^{-3/2} = 0.35$ من
القيمة القصوى.
\end{solution}

\begin{solution}{exo:b2:eruptive-styles:9}
يبلغ المدى أقصاه عند \ang{45}: $R = v^2/g = 100^2/9.81 = \qty{1.0}{km}$،
و$200^2/9.81 = \qty{4.1}{km}$. ومقاومة الهواء تقصّر الطيران، فهذان
حدّان أعليان، قريبان من الحقيقة للكتل الكبيرة وأطول بكثير من اللازم للصغيرة.
وبما أن $R = v^2/g$، فإن نطاق الخطر ينمو مع مربع سرعة الإطلاق:
فمضاعفة سرعة الانفجار تضاعف نصف قطره أربع مرات.
\end{solution}

\begin{solution}{exo:b2:eruptive-styles:10}
(أ) مع $T = T_0e^{-k_1x}$ من أجل $x < X$ و$T = T_Xe^{-k_2(x - X)}$ بعدها، حيث $T_X =
T_0e^{-k_1X}$ (الاتصال)، نستعمل $\int_0^X ye^{-ky}dy = \frac{1}{k^2} - e^{-kX}\big(\frac{X}{k} +
\frac{1}{k^2}\big)$ و$\int_X^\infty ye^{-ky}dy = e^{-kX}\big(\frac{X}{k} +
\frac{1}{k^2}\big)$:
فيعطي الجزء القريب $2T_0\big[\frac{1}{k_1^2} - e^{-k_1X}\frac{k_1X + 1}{k_1^2}\big]$
والجزء البعيد $2T_Xe^{k_2X}\int_X^\infty ye^{-k_2y}dy = 2T_X\frac{k_2X + 1}{k_2^2}$؛ ومع $T_X =
T_0e^{-k_1X}$ يكون مجموعهما صيغة إطار الطريقة.
(ب) $T_0 = \qty{0.01285}{km}$: $2T_0/k_1^2 = 0.0257/0.1332 = \qty{0.193}{km^3}$؛
$e^{-0.365\times8.56} = 0.0440$؛ $(0.138\times8.56 + 1)/0.138^2 = 114.5$
و$(0.365\times8.56 + 1)/0.365^2 = 30.96$؛ والتصحيح $0.0257\times0.0440\times83.6
= \qty{0.094}{km^3}$ و$V = \qty{0.287}{km^3}$، أي القيمة المنشورة \qty{0.29}{km^3}.
ويقع خارج $X$ المقدار $2T_X(X/k_2 + 1/k_2^2) = 2\times\num{5.65e-4}\times114.5 =
\qty{0.129}{km^3}$، أي \qty{45}{\percent}.
\end{solution}

\begin{solution}{exo:b2:eruptive-styles:11}
(أ) $100^{1/4} = 3.16$؛ ويعطي القانونان التجريبيان $100^{0.241} = 3.03$
و$100^{0.226} = 2.83$. (ب) $H \propto N^{-3/4}$: فالغلاف الجوي الأشدّ تطبّقًا
يوقف العمود أدنى؛ ومن أجل $N$ أكبر بنسبة \qty{20}{\percent}،
$1.2^{-0.75} = 0.87$، أي ارتفاع أدنى بنسبة \qty{13}{\percent}. (ج) الأعمدة الضعيفة
تحنيها الريح، فتُخفضها، في حين يحرّر الهواء الرطب حرارة كامنة
فيها، فيرفعها؛ وهذه الآثار تؤثّر على نحو مختلف في الثورانات الصغيرة والكبيرة
وليست في القانون البسيط. ثم إن المعدلات لا تُعرف إلا في حدود
عوامل من عدّة أضعاف، وهو خطأ في متغيّر المحور الأفقي ينحاز
بالميل المضبوط نحو الصفر.
\end{solution}

\begin{solution}{exo:b2:eruptive-styles:12}
(أ) $132/400\times100 = 33$ في القرن قبل 1900؛ و$95/125\times100 = 76$ في
القرن منذئذ. (ب) $33/76 = \qty{43}{\percent}$. (ج) $21/56 = \qty{38}{\percent}$:
النقص نفسه في التسجيل، بعامل اثنين إلى ثلاثة. (د) الثورانات الأكبر
تُرى من مسافة أبعد وتخلّف رواسب أثخن وأوسع تصمد للحتّ
مدة أطول؛ فثوران VEI~4 على جزيرة نائية في 1700 قد يمرّ دون تسجيل، أمّا
ثوران VEI~6 فبالكاد.
\end{solution}

\begin{solution}{pb:b2:eruptive-styles:1}
\textbf{1.} $\ln T = 5.298$، $4.605$، $3.871$، $2.639$، $1.386$. والميول بين
خطوط تساوي السمك المتتالية هي $-0.0693$ و$-0.0367$ و$-0.0308$ و$-0.0313$ لكل كيلومتر: مقطع
شديد الميل بين 10 و\qty{20}{km}، ومقطع لطيف من 40 إلى
\qty{120}{km}؛ والانكسار بين 20 و\qty{40}{km}.

\textbf{2.} $k_1 = \ln(200/100)/10 = \qty{0.0693}{km^{-1}}$؛ $T_0 = 200\,e^{0.693} =
\qty{400}{cm}$.

\textbf{3.} $167e^{-1.244} = \qty{48.1}{cm}$، $167e^{-2.488} = \qty{13.9}{cm}$،
$167e^{-3.732} = \qty{4.00}{cm}$: في حدود \qty{1}{\percent} من 48 و14 و4.0.

\textbf{4.} $400e^{-0.0693X} = 167e^{-0.0311X}$ تعطي $X = \ln(400/167)/(0.0693 -
0.0311) = 0.873/0.0382 = \qty{22.9}{km}$.

\textbf{5.} $\ln2/0.0693 = \qty{10.0}{km}$ و$\ln2/0.0311 = \qty{22.3}{km}$.

\textbf{6.} $T_0 = \qty{0.0040}{km}$: $V = 2\times0.0040/0.0693^2 =
\qty{1.67}{km^3}$.

\textbf{7.} $e^{-0.0693\times22.9} = 0.205$؛ $(0.0311\times22.9 + 1)/0.0311^2 = 1770$؛
$(0.0693\times22.9 + 1)/0.0693^2 = 539$؛ والتصحيح $2\times0.0040\times0.205\times1231 =
\qty{2.02}{km^3}$؛ $V = 1.67 + 2.02 = \qty{3.68}{km^3}$، منها للتصحيح
\qty{55}{\percent}.

\textbf{8.} $3.68/1.67 = 2.2$. فالرماد البعيد الناعم رقيق، سهل الحتّ،
يُعاد نقله أو يسقط في البحر؛ وفي الراسب القديم كثيرًا ما يضيع أو لا يُرسم
قطّ، فيُقدَّر الحجم عندئذ بأقل من قيمته بعامل اثنين أو أكثر.

\textbf{9.} يقع $\qty{3.68e9}{m^3}$ بين $10^9$ و$10^{10}$: VEI~5.

\textbf{10.} $V_{\mathrm{DRE}} = 3.68\times900/2400 = \qty{1.38}{km^3}$.

\textbf{11.} $m = \num{3.68e9}\times900 = \qty{3.31e12}{kg}$.

\textbf{12.} $M = \log_{10}(\num{3.31e12}) - 7 = 12.52 - 7 = 5.52$.

\textbf{13.} $3.68/1.3 = 2.8$ مرة.

\textbf{14.} عمود أعلى من \qty{25}{km}، وطراز بليني.

\textbf{15.} $Q = \num{3.31e12}/36000 = \qty{9.2e7}{kg/s}$؛ $I = 7.96 + 3 = 11.0$.

\textbf{16.} $0.304\times(\num{9.2e7})^{0.241} = 0.304\times83.0 = \qty{25.2}{km}$.

\textbf{17.} $0.345\times(\num{9.2e7})^{0.226} = 0.345\times63.1 = \qty{21.8}{km}$.

\textbf{18.} الارتفاع المقيس \qty{31}{km} أعلى بنسبة \qty{23}{\percent} من التنبّؤ
الأول وبنسبة \qty{42}{\percent} من الثاني، في حدود التشتّت
\qtyrange{53}{60}{\percent} للقانونين: فهما متّسقان.

\textbf{19.} $Q = (31/0.304)^{1/0.241} = 102.0^{4.149} = \qty{2.2e8}{kg/s}$، أي 2.3 ضعف
المعدل المتوسط. فالقطع تسجّل أعلى عمود، عند ذروة
الثوران، في حين يأخذ السؤال~15 المتوسط على الطور كله؛ وقد يكون جزء من الكتلة
مفقودًا من الراسب (الرماد البعيد، التدفّقات الفتاتية البركانية).

\textbf{20.} نعم: من 22 إلى \qty{31}{km}، أعلى بكثير من التروبوبوز عند \qty{16}{km}.
والغاز المحمول إلى الستراتوسفير ينتشر حول الكرة الأرضية؛ والغاز الأساسي هو
ثاني أكسيد الكبريت، والعشرون مليون طن التي حقنها بيناتوبو هناك في 1991
بردت الكرة الأرضية بنحو نصف درجة مدة سنتين.

\textbf{21.} $\lambda_3 = 399/75 = 5.32$، $\lambda_4 = 95/125 = 0.760$، $\lambda_5 = 22/225 =
0.0978$، $\lambda_6 = 8/525 = 0.0152$ في السنة؛ ولوغاريتماتها $0.726$، $-0.119$، $-1.010$،
$-1.817$.

\textbf{22.} يعطي المستقيم $10^{0.73} = 5.37$ و$10^{-0.12} = 0.759$ و$10^{-0.97} =
0.107$ و$10^{-1.82} = 0.0151$: في حدود \qty{1}{\percent} و\qty{0}{\percent}
و\qty{10}{\percent} و\qty{1}{\percent} من الأعداد.

\textbf{23.} يبقى سجلّ الثورانات الكبيرة مكتملًا أبعد في الماضي،
لأنها تُلاحظ من مسافة أبعد وتصمد رواسبها مدة أطول؛ فنافذة
تبدأ في 1500 تُنقص عدّ ثورانات VEI~3 كثيرًا، في حين أن نافذة
تبدأ في 1950 لا تكاد تضمّ أي ثوران VEI~6.

\textbf{24.} $\lambda_5 = \num{0.107}$ في السنة، أي فترة عودة قدرها $1/0.107 =
\num{9.3}\text{ سنة}$؛ و$1 - e^{-1.07} = 1 - 0.343 = \qty{66}{\percent}$ في عشر سنوات.

\textbf{25.} $\lambda_7 = 10^{3.28 - 5.95} = 10^{-2.67} = 0.00214$ في السنة، أي ثوران كل
470 سنة؛ ويعطي السجلّ $2/1025 = 0.00195$ في السنة: اتفاق في حدود
\qty{10}{\percent}، لنقطة لم تُستعمل في الضبط.

\textbf{26.} $\lambda_6 = 10^{3.28 - 5.10} = 10^{-1.82} = 0.0151$ في السنة؛ $\lambda t = 1.51$
و$1 - e^{-1.51} = 1 - 0.220 = 0.780$: فاحتمال ثوران واحد على الأقل
من VEI~6 أو أكثر في مكان ما على الأرض خلال المئة سنة المقبلة هو
\textbf{\qty[text-series-to-math=true]{78}{\percent}}. ويستند إلى معدل ثابت،
وإلى ثورانات مستقلة بعضها عن بعض (عملية بواسون)، وإلى
فهرس مكتمل ضمن نوافذه.
\end{solution}
